\documentclass[10pt,letterpaper]{article}
\usepackage[utf8]{inputenc}
\usepackage[T1]{fontenc}
\usepackage[margin=0.65in]{geometry}
\usepackage{amsmath,amssymb}
\usepackage{booktabs}
\usepackage{hyperref}

\title{\textbf{NeuroSpace Lab: 3D Neural Network Topology, Structural Plasticity, and Six-Layered Cerebral Cortex Simulation in Rust}}
\author{Bryan P. White (ORCID: \href{https://orcid.org/0000-0003-4191-3511}{0000-0003-4191-3511}) \\ CRV Science Labs \\ \small CRV Science Apps Hub: \url{https://crvscience.com/apps} \\ \small DOI: \href{https://doi.org/10.5281/zenodo.24263468}{10.5281/zenodo.24263468}}
\date{November 2026 \\ \small Preprint ID: APP-PREPRINT-19 (v1.6.0 Software Preprint) \\ \small Categories: q-bio.NC \textbullet{} cs.NE \textbullet{} cs.GR}

\begin{document}
\maketitle

\begin{abstract}
NeuroSpace Lab (v1.6.0) is an open-source 3D neural network topology and cortical wetware workstation implemented in Rust to unify connectionist architectures and biophysical cortical simulation inside a shared interactive coordinate system. Spanning 15 architectures (M0--M14) grouped across five evolutionary tiers---from direct-gate perceptrons, residual CNNs, and GQA sparse Mixture-of-Experts transformers (M14) to STDP spiking nets, liquid time-constant ODEs, predictive coding, HTM columns, connectomic organoids, and six-layered Cerebral Cortex Grey Matter (M13, L1--L6 laminae with Basket/Martinotti inhibition and astrocyte modulation)---the platform evaluates every model across seven deterministic benchmark domains (D1--D7) alongside live synaptic sprouting, pruning, neurogenesis, ChR2/NpHR optogenetic clamping, and causal lesions. On 80-step CLI sweeps (seed = 2026) of D4 FloraMorph botanical classification, Hierarchical Predictive Coding M10 reached 50.00\% accuracy and Connectomic Organoid M12 reached 45.83\% (followed by M0/M5 at 35.42\% and Frontier LLM M14 at 29.17\%), whereas holding M13 laminar cortex fixed yielded 66.30\% R2 on D5 Willamette regression, 69.21\% score on D6 Robotics LiDAR-12, and 59.80\% concordance on D7 ChemGraph QSAR under GPL-3.0-or-later licensing.
\end{abstract}

\section{Introduction}
Modern neural computation is commonly taught and evaluated inside disconnected silos. Classical machine-learning frameworks treat neural networks as abstract tensor pipelines tuned purely for loss minimization, whereas computational neuroscience toolkits emphasize compartmental biophysics without exposing comparable end-to-end task performance across vision, regression, autonomous control, and molecular property prediction. When students or researchers move from a multi-layer perceptron or sparse-attention transformer [4] to spike-timing-dependent plasticity (STDP) [2] , liquid time-constant ordinary differential equations [3] , predictive coding [1] , or Thousand-Brains neocortical columns [5] , they rarely observe those architectures competing side by side on identical deterministic feature streams inside the same three-dimensional spatial coordinate system. NeuroSpace Lab (v1.6.0) bridges this conceptual divide in a unified interactive 3D simulation environment. Developed within the CRV Science Open-Source Apps ecosystem [6--25] , NeuroSpace Lab places 15 neural architectures ( M0 -- M14 ) spanning five evolutionary tiers inside a shared topological arena backed by seven deterministic scientific benchmark domains ( D1 -- D7 ), real-time structural plasticity (synaptic sprouting, pruning, and neurogenesis), causal perturbation tools, and reproducible experiment checkpoints. Rather than treating visualization as a decorative afterthought, NeuroSpace Lab couples every learning hyperparameter and biophysical control---learning rate, momentum, pruning threshold, neurogenesis rate, cortical capacity, and adversarial perturbation radius---directly to 3D soma geometry, dendritic morphology, and excitatory/inhibitory axonal routing. Investigators can hold a benchmark domain fixed to compare all 15 architectures head-to-head, or hold an architecture such as six-layered Cerebral Cortex Grey Matter ( M13 ) fixed to study generalization across all seven scientific domains.

\section{Architecture \& Methods}
NeuroSpace Lab organizes its 15 models into five conceptual tiers that trace the progression from classical connectionist layers to biologically grounded cortical circuits and modern sparse foundation models (Table 1). On every simulation step, an active model receives a deterministic feature stream, computes task-appropriate error signals---Cross-Entropy for multi-class classification, Huber/RMSE loss for hedonic regression, control cost J ( u ) for kinematic navigation, and multi-task MAE for cheminformatics---and updates its synaptic weights and structural connectivity according to its native plasticity rule. Version 1.6.0 introduces two conceptual advances for anatomical fidelity and visual interpretability. First, Calm, Sparse Macroscopic Flow Sampling replaces dense per-synapse particle barrages with smooth, sine-faded directional flow wisps sampled across a representative subset of pathways with refractory rest intervals. By limiting concurrent tracers, macroscopic feedforward sweeps and recurrent feedback volleys remain immediately legible without obscuring individual somata. Second, Tier V ( M13 ) Cerebral Cortex Grey Matter models a canonical six-layered neocortical column extending from Pia Mater ( +Y ) down to subcortical White Matter ( -Y ): thalamic sensory afferents synapse onto Layer 4 ( L4 ) Spiny Stellate Granules, ascend to L2/3 Association Pyramidals, project to L5 Thick-Tufted Pyramidals exhibiting apical Ca 2+ burst signaling, and close corticothalamic loops via L6 Fusiforms alongside L1 horizontal apical tufts, GABAergic Basket and Martinotti interneurons, and protoplasmic astrocytes.

\section{Capabilities \& Benchmark Results}
Deterministic headless CLI sweeps ( seed = 2026 , 80 training steps per configuration) demonstrate how architectural inductive bias interacts with task geometry. Under locked dataset D4 (FloraMorph PNW 10-feature botanical taxonomy across 6 classes), Hierarchical Predictive Coding ( M10 ) reaches 50.00\% validation accuracy and Connectomic Organoid-Astrocyte networks ( M12 ) reach 45.83\% within 80 steps, followed by compact direct-gate perceptrons ( M0 , 35.42\% ), PGD-robust networks ( M5 , 35.42\% ), and Frontier Autoregressive LLM with GQA + Sparse MoE ( M14 , 29.17\% ). Conversely, locking Tier V Cerebral Cortex Grey Matter ( M13 ) and sweeping all seven datasets ( D1 -- D7 ) highlights the strong advantage of laminar feedback and apical burst integration on continuous and multi-output scientific control/regression problems. On D5 Willamette Hedonic home valuation, M13 reaches 66.30\% R 2 fit (Huber/RMSE loss 0.1198 ); on D6 Robotics LiDAR-12 autonomous navigation, M13 achieves 69.21\% collision-free score (control cost 0.1835 ); and on D7 ChemGraph QSAR multi-property prediction (BBB permeability, HOMO-LUMO gap, aqueous solubility), M13 reaches 59.80\% concordance (multi-task MAE 0.2181 ).

\section{Causal Lesions, Optogenetics \& Interface Gallery}
Beyond forward benchmarking, NeuroSpace Lab embeds an interactive causal perturbation console for testing circuit necessity and sufficiency. Users can switch between soma inspection, physical ablation, Channelrhodopsin-2 (ChR2) optical excitation, and Halorhodopsin (NpHR) hyperpolarizing silencing directly on 3D somata, or trigger macro-level interventions---severing top-down feedback trunks, cutting lateral skip connections, or inducing a 25\% focal cortical stroke---followed by single-click tissue recovery. Comparing performance degradation before and after severing feedback loops makes concepts such as predictive error correction [1] and cortical column voting [5] immediately tangible.

\section{Discussion}
Evaluating all fifteen neural architectures inside identical deterministic feature streams exposes an informative inductive-bias trade-off between early parameter efficiency and hierarchical cortical depth. On fixed 80-step sweeps of D4 FloraMorph botanical classification, Hierarchical Predictive Coding M10 ( 50.00\% ), Connectomic Organoid M12 ( 45.83\% ), and compact M0 / M5 ( 35.42\% ) converge rapidly, whereas the deeper six-layered Cerebral Cortex microcircuit M13 (61 neurons / 425 synapses) attains 12.50\% within the same 80-step budget on D4 ---reflecting slower early optimization across feedforward-to-feedback pathways ( L4 → L2/3 → L5 → L6 ). Conversely, on higher-dimensional continuous control and multi-task regression domains such as D5 Willamette valuation ( 66.30\% R 2 ), D6 LiDAR-12 robotics ( 69.21\% ), and D7 ChemGraph QSAR ( 59.80\% concordance), laminar recurrent feedback and canonical cortical integration support strong multi-output generalization. Equally important, pairing scalar benchmark accuracy with interactive Channelrhodopsin-2 stimulation, Halorhodopsin silencing, and structural cortical stroke lesions clarifies circuit mechanisms obscured by loss curves alone: severing top-down apical feedback trunks degrades downstream motor and regression readouts selectively across layers rather than uniformly. Current evaluations are intentionally bounded to fixed 80-step deterministic sweeps and simplified point/compartmental soma dynamics on curated scientific fixtures rather than full Hodgkin-Huxley ion-channel conductances; extending training horizons, incorporating dendritic calcium spikes, and benchmarking online structural plasticity across shifting non-stationary distributions remain key directions for future releases.

\section{Software Availability \& Network Integration}
NeuroSpace Lab (v1.6.0) is released in full Rust source under the GNU General Public License v3.0 or later ( GPL-3.0-or-later ), supporting deterministic seeded sweeps, portable local experiment snapshots, and standalone offline execution. Within the broader twenty-program CRV Science Open-Source Apps suite at https://crvscience.com/apps [6--25] , NeuroSpace Lab complements molecular and biophysical simulators such as BeoFold, SAR-Atlas, and Causal-Atlas, mathematical and bioacoustic workbenches such as Calc-Lab, OSABT, PopBio-Studio, and SonoWing, as well as systems and multi-agent simulators including HelioProbe, QuantaCrypt-Studio, TerraTwin, Urban-Grid, SwarmForge, and WildFront BattleGrounds.

Four-Pillar Engineering Refinement \& Verification: Non-blocking incremental batch sweeps, camera preset/tool hotkeys, and right-drag pan pair with true gradient-aligned FGSM/PGD adversarial evaluation, wired immunity and free-energy goals, hop-staged backprop, exponential LIF decay with trace STDP, L1 apical feedback, 3D Bezier/pyramidal geometry with frustum clipping, lesion tags, Markdown guides, and 20/20 passing tests.

\begin{thebibliography}{28}
\bibitem{ref1} Rao, R. P. N., \& Ballard, D. H. (1999). Predictive coding in the visual cortex: a functional interpretation of some extra-classical receptive-field effects. Nature Neuroscience, 2(1), 79--87 \url{https://doi.org/10.1038/4580}

\bibitem{ref2} Bi, G.-q., \& Poo, M.-m. (1998). Synaptic Modifications in Cultured Hippocampal Neurons: Dependence on Spike Timing, Synaptic Strength, and Postsynaptic Cell Type. Journal of Neuroscience, 18(24), 10464--10472 \url{https://doi.org/10.1523/JNEUROSCI.18-24-10464.1998}

\bibitem{ref3} Hasani, R., Lechner, M., Amini, A., Rus, D., \& Grosu, R. (2021). Liquid Time-constant Networks. Proceedings of the AAAI Conference on Artificial Intelligence, 35(9), 7657--7666 \url{https://doi.org/10.1609/aaai.v35i9.16936}

\bibitem{ref4} Vaswani, A., et al. (2017). Attention Is All You Need. Advances in Neural Information Processing Systems (NeurIPS 2017) \url{https://doi.org/10.48550/arXiv.1706.03762}

\bibitem{ref5} Hawkins, J., Ahmad, S., \& Cui, Y. (2017). A Theory of How Columns in the Neocortex Enable Learning the Structure of the World. Frontiers in Neural Circuits, 11, 81 \url{https://doi.org/10.3389/fncir.2017.00081}

\bibitem{ref6} White, B. P. (2026). The CRV Science App Network: Twenty Open-Source Rust Programs for Science, Data, and Simulation. CRV Science App Preprint APP-PREPRINT-01 [cs.SE, cs.CE, q-bio.QM], v1.1.0. Hosted via CRV Science Apps Hub (https://crvscience.com/apps) \url{https://doi.org/10.5281/zenodo.24222398}

\bibitem{ref7} White, B. P. (2026). BeoFold: Simulating How Proteins Fold Using Distributed Computing and State Networks. CRV Science App Preprint APP-PREPRINT-02 [q-bio.BM, physics.chem-ph, cs.DC], v1.4.0. Hosted via CRV Science Apps Hub (https://crvscience.com/apps) \url{https://doi.org/10.5281/zenodo.24222370}

\bibitem{ref8} White, B. P. (2026). Calc-Lab: Interactive Calculus, Linear Algebra, and Differential Equation Solver Workbench. CRV Science App Preprint APP-PREPRINT-03 [math.NA, cs.MS, cs.CY], v0.5.0. Hosted via CRV Science Apps Hub (https://crvscience.com/apps) \url{https://doi.org/10.5281/zenodo.24222372}

\bibitem{ref9} White, B. P. (2026). Causal-Atlas: Tracing How Genes and Drugs Turn Biological Pathways On and Off. CRV Science App Preprint APP-PREPRINT-04 [q-bio.MN, q-bio.GN, stat.ML], v0.2.0. Hosted via CRV Science Apps Hub (https://crvscience.com/apps) \url{https://doi.org/10.5281/zenodo.24222376}

\bibitem{ref10} White, B. P. (2026). Cit-Atlas: Mapping Citation Networks and Research Impact Across Scientific Papers. CRV Science App Preprint APP-PREPRINT-05 [cs.DL, cs.SI, physics.soc-ph], v1.8.0. Hosted via CRV Science Apps Hub (https://crvscience.com/apps) \url{https://doi.org/10.5281/zenodo.24222380}

\bibitem{ref11} White, B. P. (2026). EpiLedger: Privacy-Safe Disease Tracking, Clinical Text Screening, and Outbreak Monitoring. CRV Science App Preprint APP-PREPRINT-06 [cs.CR, q-bio.PE, stat.AP], v2.0.0. Hosted via CRV Science Apps Hub (https://crvscience.com/apps) \url{https://doi.org/10.5281/zenodo.24222382}

\bibitem{ref12} White, B. P. (2026). ITME: Modeling Evolutionary Regimes, Contingency, and Substrate Complexity Across Molecular, Hominid, and AI Worlds. CRV Science App Preprint APP-PREPRINT-07 [q-bio.PE, nlin.AO, physics.bio-ph], v2.3.0. Hosted via CRV Science Apps Hub (https://crvscience.com/apps) \url{https://doi.org/10.5281/zenodo.24222386}

\bibitem{ref13} White, B. P. (2026). OSABT: A 20-Module Biology and Biostatistics Workbench for Molecular Labs, Genetics, Kinetics, and Ecology. CRV Science App Preprint APP-PREPRINT-08 [q-bio.QM, stat.AP, cs.SE], v2.0.0. Hosted via CRV Science Apps Hub (https://crvscience.com/apps) \url{https://doi.org/10.5281/zenodo.24222388}

\bibitem{ref14} White, B. P. (2026). SAR-Atlas: Drug Design, Body Distribution Modeling, and Candidate Screening in Rust. CRV Science App Preprint APP-PREPRINT-09 [physics.chem-ph, q-bio.BM, cs.LG], v3.8.0. Hosted via CRV Science Apps Hub (https://crvscience.com/apps) \url{https://doi.org/10.5281/zenodo.24222390}

\bibitem{ref15} White, B. P. (2026). TerraTwin: Fast Climate Projections for Global Warming, Sea-Level Rise, and Arctic Change. CRV Science App Preprint APP-PREPRINT-10 [physics.ao-ph, cs.CE, stat.AP], v2.5.0. Hosted via CRV Science Apps Hub (https://crvscience.com/apps) \url{https://doi.org/10.5281/zenodo.24222392}

\bibitem{ref16} White, B. P. (2026). Urban-Grid: Modeling City Traffic, Emergency Response Coverage, and Earthquake Evacuation. CRV Science App Preprint APP-PREPRINT-11 [physics.soc-ph, cs.MA, cs.CE], v5.2.0. Hosted via CRV Science Apps Hub (https://crvscience.com/apps) \url{https://doi.org/10.5281/zenodo.24222394}

\bibitem{ref17} White, B. P. (2026). WildFront BattleGrounds: Turn-Based Hex Strategy Game with Voxel Graphics and Procedural Map Generation. CRV Science App Preprint APP-PREPRINT-12 [cs.GT, cs.GR, math.CO], v3.4.1. Hosted via CRV Science Apps Hub (https://crvscience.com/apps) \url{https://doi.org/10.5281/zenodo.24222396}

\bibitem{ref18} White, B. P. (2026). AeroFlyway: Tracking Bird Migration with Weather Radar, Skyglow Maps, and Nighttime Flight Calls. CRV Science App Preprint APP-PREPRINT-13 [q-bio.PE, physics.ao-ph, cs.GR], v0.6.0. Hosted via CRV Science Apps Hub (https://crvscience.com/apps) \url{https://doi.org/10.5281/zenodo.24222366}

\bibitem{ref19} White, B. P. (2026). HelioProbe: Simulating Spacecraft Orbits, Deep-Space Telescopes, and Comet Tails in 3D. CRV Science App Preprint APP-PREPRINT-14 [astro-ph.EP, astro-ph.IM, cs.GR], v2.3.0. Hosted via CRV Science Apps Hub (https://crvscience.com/apps) \url{https://doi.org/10.5281/zenodo.24244083}

\bibitem{ref20} White, B. P. (2026). PopBio-Studio: Interactive Population Biology, Adaptive Landscapes, and Genetic Skyline Simulator. CRV Science App Preprint APP-PREPRINT-15 [q-bio.PE, q-bio.QM, stat.AP], v1.2.0. Hosted via CRV Science Apps Hub (https://crvscience.com/apps) \url{https://doi.org/10.5281/zenodo.24244085}

\bibitem{ref21} White, B. P. (2026). QuantaCrypt-Studio: Simulating Quantum Circuits, Quantum Key Distribution, and Post-Quantum Cryptography. CRV Science App Preprint APP-PREPRINT-16 [quant-ph, cs.CR, cs.MS], v2.2.0. Hosted via CRV Science Apps Hub (https://crvscience.com/apps) \url{https://doi.org/10.5281/zenodo.24244089}

\bibitem{ref22} White, B. P. (2026). SonoWing: Analyzing Bat Echolocation and Bird Calls with Spectrograms and Harmonic Tracking. CRV Science App Preprint APP-PREPRINT-17 [q-bio.NC, q-bio.PE, cs.SD], v1.1.0. Hosted via CRV Science Apps Hub (https://crvscience.com/apps) \url{https://doi.org/10.5281/zenodo.24244091}

\bibitem{ref23} White, B. P. (2026). SwarmForge: 3D Isometric Automated Factory and Swarm Robotics Simulation Game. CRV Science App Preprint APP-PREPRINT-18 [cs.MA, cs.RO, cs.GT], v1.3.0. Hosted via CRV Science Apps Hub (https://crvscience.com/apps) \url{https://doi.org/10.5281/zenodo.24244093}

\bibitem{ref24} White, B. P. (2026). TrophicHearth: Nine-Biome 3D Living Trophic Diorama, Fire-Pit Hearth, and Bioacoustic Ecosystem Simulator in Rust. CRV Science App Preprint APP-PREPRINT-20 [q-bio.PE, cs.GR, cs.MA], v0.9.3. Hosted via CRV Science Apps Hub (https://crvscience.com/apps) \url{https://doi.org/10.5281/zenodo.24359107}

\bibitem{ref25} White, B. P. (2026). ConnectomeAtlas: Standalone 3D Multi-Scale Brain Connectome, Persistent Homology, and Lesion Workstation in Rust. CRV Science App Preprint APP-PREPRINT-21 [q-bio.NC, cs.SI, cs.GR], v0.8.0. Hosted via CRV Science Apps Hub (https://crvscience.com/apps) \url{https://doi.org/10.5281/zenodo.24359109}

\bibitem{ref26} White, B. P. (2026). PelagicGlobe: 3D Ocean Currents Globe, Conformal Spilhaus World Ocean Projection, and NOAA ENSO Workstation in Rust. CRV Science App Preprint APP-PREPRINT-22 [physics.ao-ph, cs.GR, cs.CE], v1.1.0. Hosted via CRV Science Apps Hub (https://crvscience.com/apps) \url{https://doi.org/10.5281/zenodo.24442089}

\bibitem{ref27} White, B. P. (2026). CanopyForge: 3D Isometric Pacific Northwest Old-Growth Watershed Recovery and Ecological Stewardship Simulator in Rust. CRV Science App Preprint APP-PREPRINT-23 [q-bio.PE, cs.GR, cs.MA], v1.0.0. Hosted via CRV Science Apps Hub (https://crvscience.com/apps) \url{https://doi.org/10.5281/zenodo.24442091}

\bibitem{ref28} White, B. P. (2026). GenBankVault: Safe-Rust NCBI Five-Gate Genomic Downloader, Multi-FASTA Set Organizer, and Offline SQLite Synchronization Workstation. CRV Science App Preprint APP-PREPRINT-24 [q-bio.GN, cs.DB, cs.GR], v1.1.0. Hosted via CRV Science Apps Hub (https://crvscience.com/apps) \url{https://doi.org/10.5281/zenodo.24451850}
\end{thebibliography}
\end{document}
